\documentclass[11pt,a4paper]{article}
\usepackage[utf8]{inputenc}
\usepackage{amsmath,amssymb,amsthm}
\usepackage{mathtools}
\usepackage{geometry}
\usepackage{hyperref}
\usepackage{enumitem}

\geometry{margin=2.5cm}

\newtheorem{theorem}{Theorem}[section]
\newtheorem{lemma}[theorem]{Lemma}
\newtheorem{proposition}[theorem]{Proposition}
\newtheorem{definition}[theorem]{Definition}
\newtheorem{remark}[theorem]{Remark}

\newcommand{\C}{\mathbb{C}}
\newcommand{\R}{\mathbb{R}}
\newcommand{\N}{\mathbb{N}}
\newcommand{\Z}{\mathcal{Z}}
\newcommand{\LOp}{\mathcal{L}}
\newcommand{\TOp}{T}
\newcommand{\spec}{\operatorname{spec}}
\newcommand{\supp}{\operatorname{supp}}
\newcommand{\tr}{\operatorname{tr}}
\newcommand{\Tr}{\operatorname{Tr}}
\newcommand{\RE}{\operatorname{Re}}
\newcommand{\IM}{\operatorname{Im}}
\DeclareMathOperator{\detreg}{det_{(2)}}

\title{\bf On the Possibility of a Variational Mechanism\\ in the Iota Walk Programme}
\author{}
\date{\today}

\begin{document}
\maketitle

\begin{abstract}
We analyse the proposed connection between the greedy condition (GC) on the iota walk and the spectral radius of the associated transfer operator \(L_s\) from the HST framework.  We show that any operator representation of the iota terms of the form \(v_n(s) = \ell(L_s^{\,n-1} v_0)\) with a \textbf{continuous} linear functional \(\ell\) on the Banach space on which \(L_s\) is bounded forces \(\rho(L_s) \ge 1\) for \emph{all} \(\sigma > 0\), contradicting the well‑known strict contraction property of \(L_s\) for \(\sigma > 1/2\).  Hence such a representation cannot exist.  The actual relationship between the iota walk and the transfer operator must involve unbounded functionals, which prevents a direct translation of the greedy condition into a bound on the spectral radius.  We discuss the implications for the research programme and outline alternative approaches.
\end{abstract}

\section{Introduction}
The goal of the iota walk programme is to prove that if \(s = \sigma + it\) is a zero of \(\eta(s) = \sum_{n=1}^\infty (-1)^{n-1} n^{-s}\) and the greedy condition (GC)
\[
\langle \mathbf{S}_{N-1}(s),\, v_N(s)\rangle \le 0 \qquad \forall N\ge 1
\]
holds, then necessarily \(\sigma = 1/2\).  A central part of the programme has been to connect the iota walk to a transfer operator \(L_s\) constructed in the Harmonic Sine Transform (HST) framework \cite{HST}, with the hope of converting GC into a spectral constraint on \(L_s\).  In earlier notes, Task~1 \cite{Task1} claimed an operator representation of the iota terms as matrix elements of powers of \(L_s\) via a continuous linear functional.  The present note demonstrates that such a representation is impossible if the functional is continuous.  This invalidates the simple variational mechanism and forces a re‑evaluation of the strategy.

\section{The Transfer Operator \(L_s\) and its Spectral Radius}
The HST paper constructs, for \(\RE s > 1/2\), a trace‑class operator \(L_s\) on the Hardy space \(H^2(\mathbb{D})\) whose Fredholm determinant satisfies
\[
\det(I - L_s) = \frac{\zeta(s)}{\zeta(2s)} \, e^{P(s)},
\tag{1}
\]
with \(P(s)\) entire and non‑vanishing on the critical line.  The operator \(L_s\) depends analytically on \(s\) and, for real \(\sigma\), enjoys the following properties:
\begin{itemize}
\item For \(\sigma > 1/2\), \(L_s\) is a strict contraction: \(\rho(L_s) < 1\).
\item For \(\sigma = 1/2\), \(L_{1/2+it}\) is similar to a unitary operator; in particular, \(\rho(L_{1/2+it}) = 1\).
\item For \(\sigma < 1/2\), the operator is not bounded on the same space; its spectral radius exceeds \(1\) on suitable Banach spaces.
\end{itemize}
These facts are standard for transfer operators of the Gauss map and carry over to the HST variant \cite{Mayer,Efrat}.

\section{The Iota Walk and its Terms}
The Dirichlet eta function is \(\eta(s) = (1-2^{1-s})\zeta(s)\).  For \(\sigma > 0\), it is represented by the convergent alternating series
\[
\eta(s) = \sum_{n=1}^\infty (-1)^{n-1} n^{-s}.
\]
The iota walk is the sequence of partial sums
\[
\mathbf{S}_N(s) = \sum_{n=1}^{N} v_n(s), \qquad
v_n(s) = (-1)^{n-1} n^{-s}.
\]

\section{An Attempted Operator Representation}
Task~1 \cite{Task1} proposed the existence of a vector \(v_0 \in H^2(\mathbb{D})\) and a continuous linear functional \(\ell\) on that space such that
\[
\ell(L_s^{\,n-1} v_0) = v_n(s) \qquad (n\ge 1).
\tag{2}
\]
If this representation were true, then for every \(s\) with \(\sigma > 0\), the power series \(\sum_{n=1}^\infty v_n(s) z^{n-1}\) would be given by \(\ell( (I - z L_s)^{-1} v_0 )\).  The radius of convergence of that power series is exactly \(1\), because \(v_n(s) \sim n^{-\sigma}\) up to oscillation, which is not summable for \(|z|>1\).  Hence the operator \((I - z L_s)\) must be invertible for all \(|z| < 1\) but fails to be invertible for \(|z| > 1\).  Consequently, the spectral radius of \(L_s\) must satisfy \(\rho(L_s) \ge 1\).  This conclusion holds for \emph{all} \(\sigma > 0\) for which the representation (2) is valid.

However, it is a mathematically established fact that for \(\sigma > 1/2\), the operator \(L_s\) on \(H^2(\mathbb{D})\) has spectral radius strictly less than \(1\).  Hence the representation (2) cannot hold with a \textbf{continuous} functional \(\ell\) on a space where \(L_s\) is a bounded operator.  Any correct operator connection must involve an unbounded linear functional, or must be formulated on a space where \(L_s\) is not bounded for \(\sigma > 1/2\).

\section{Implications for the Variational Mechanism}
The impossibility of a continuous representation fatally undermines the elementary variational argument that GC would bound the spectral radius.  The sequence \(v_n(s)\) decays polynomially, not exponentially, for all \(\sigma > 0\); it is therefore completely insensitive to the exponential contraction or expansion of the operator norm.  Any attempt to deduce \(\rho(L_s) \le 1\) from GC using a direct spectral radius estimate on the same space cannot succeed.

One might hope to place \(L_s\) on a larger space where it is bounded for \(\sigma \le 1/2\) and where the evaluation functional becomes continuous.  For example, on a space of sequences with a weighted sup‑norm, the transfer operator could be represented by an infinite matrix, and the term \(v_n(s)\) could be a coordinate of \(L_s^{\,n-1} v_0\).  In such a space, however, the spectral radius of \(L_s\) is likely \(+\infty\) for \(\sigma\) small, and the relationship (1) with \(\zeta(s)\) is lost.  The link between the iota walk and the zeta function would then no longer be available through the operator \(L_s\).

\section{A Possible Exit: Projective Metrics on Cones}
A more refined approach, outlined in earlier notes, uses Hilbert's projective metric on a cone preserved by \(L_s\).  For real \(\sigma\) the operator \(L_\sigma\) (with \(t=0\)) is a positive operator on a suitable function space, and its projective contraction coefficient is strictly monotone in \(\sigma\), passing through \(1\) exactly at \(\sigma = 1/2\).  For complex \(t\), the operator is no longer positive, but its action on the cone may be replaced by the action of a related self‑adjoint operator on the critical line.  The greedy condition, translated into a condition on the orbit of a specific vector within the cone, might still enforce a projective non‑expansion, which at \(\sigma = 1/2\) is the borderline case.  However, the translation of GC into a cone condition remains entirely conjectural.

\section{Conclusion}
We have shown that a simple continuous operator representation of the iota terms cannot hold without violating known spectral properties of the HST transfer operator.  This negates the elementary variational argument and forces the programme to either abandon the direct operator connection or to develop a much more subtle relationship, perhaps based on projective geometry or unbounded operator theory.  Until such a relationship is rigorously established, the claimed connection between the greedy condition and the spectral radius of \(L_s\) remains unsupported, and the iota walk programme does not provide a proof of the Riemann Hypothesis.

\begin{thebibliography}{9}
\bibitem{Task1} V.\ Geere, \emph{Task 1 -- Explicit Connection Between the Iota Walk and the Transfer Operator}, research note, 2026.
\bibitem{HST} V.\ Geere, \emph{Harmonic Sine Transform: a categorical Hilbert--Pólya approach to the Riemann Hypothesis}, preprint, 2026.
\bibitem{Mayer} D.\ Mayer, \emph{Continued fractions and related transformations}, in: \textit{Ergodic Theory, Symbolic Dynamics, and Hyperbolic Spaces}, Oxford Univ.\ Press, 1991.
\bibitem{Efrat} I.\ Efrat, \emph{Dynamics of the continued fraction map and the spectral theory of the transfer operator}, Nonlinearity \textbf{6} (1993), 619--634.
\end{thebibliography}

\end{document}